\documentclass[a4paper,addpoints,14pt]{exam}
%\documentclass[answers,addpoints,14pt]{exam}
\unframedsolutions

\usepackage{array,amsmath,fancybox,amsfonts,graphicx,amsthm,amssymb,wasysym}
\usepackage{latexsym,slashed,subfigure,calc}
\usepackage[spanish]{babel}
\usepackage[utf8]{inputenc}
%\usepackage{fancyhdr,float,hyperref}
%\usepackage{lettrine}

\usepackage{enumitem,mdframed}

%entorno para química: texto y math mode
%\newcommand*\chem[1]{\textup{#1}}
\newcommand*\chem[1]{\ensuremath{\mathrm{#1}}}

%tamaños fuentes extra: 8,9,14,17,20
	\usepackage{extsizes}

%display style in long fractions
\newcommand\ddfrac[2]{\frac{\displaystyle #1}{\displaystyle #2}}

%tabla generada
\usepackage[table,xcdraw]{xcolor}

%lewis y química
\usepackage{chemformula,elements}
%chemformula, provides \ch{} with the !(<below>)(<formula>) syntax and \chlewis;
%elements, provides \elconf and \writeelconf.
%ex:
%\ch{
%  !(\elconf{Li})( "\chlewis{180.}{Li}" ) +
%  !(\elconf{F})( "\chlewis{0.90:180:270:}{F}" )
%   ->
%  !(\writeelconf{2})( Li+ ) +
%  !(\writeelconf{2,2+6})( "\chlewis{0:90:180:270:}{F}" {}- )
%}

%multicol
\usepackage{multicol}
  

%símbolo bonito para qed
\usepackage{amsthm}
	%\newcommand{\mc}[3]{\multicolumn{#1}{#2}{#3}}
\newcommand{\halmos}{\vspace*{-6mm}
	
					\begin{flushright}
						\rule{1mm}{2.5mm} 
					\end{flushright}
					\vspace*{-4mm}
						
					  }
\renewcommand{\qed}{\halmos}

 
%página blanco
\usepackage{afterpage}
%uso \afterpage{\blankpage}

\newcommand\blankpage{%
    \null
    \thispagestyle{empty}%
    \addtocounter{page}{0}%
    \newpage}





 %caja bonita
  \usepackage[many]{tcolorbox}
\definecolor{titlebg}{RGB}{80,100,72}
\newtcolorbox{learn}{
  breakable,
  enhanced,
  arc=0pt,
  outer arc=0pt,
  colframe=titlebg,
  colback=titlebg!7,
  overlay unbroken and first={
    \node[
      draw=titlebg,
      fill=titlebg,
      rotate=270,
      anchor=north west,
      text=white,
      font=\bfseries
    ]
    at (frame.north west)  
    {DATOS};
  }
}


\newtcolorbox{formulacion}{
  breakable,
  enhanced,
  arc=0pt,
  outer arc=0pt,
  colframe=titlebg,
  colback=titlebg!7,
  overlay unbroken and first={
    \node[
      draw=titlebg,
      fill=titlebg,
      rotate=270,
      anchor=north west,
      text=white,
      font=
    ]
    at (frame.north west)  
    {\hspace*{3ex}};
  }
}

\newtcolorbox{problemas}{
  breakable,
  enhanced,
  arc=0pt,
  outer arc=0pt,
  colframe=titlebg,
  colback=titlebg!7,
  overlay unbroken and first={
    \node[
      draw=titlebg,
      fill=titlebg,
      rotate=270,
      anchor=north west,
      text=white,
      font=
    ]
    at (frame.north west)  
    {AVISO};
  }
}
	\usepackage{anysize}
	%\marginsize{left}{right}{top}{bottom}
    \marginsize{.9cm}{.8cm}{1cm}{1cm}
    
\usepackage[spanish]{babel}
\usepackage[utf8]{inputenc}

\usepackage{sidecap}
\usepackage{mwe,float,afterpage}

%redefiniciones
\pointpoints{punto}{puntos}
\bonuspointpoints{punto}{puntos}
\hqword{Pregunta} %substitutes text for \Question"
%\vpgword{Pág.}%substitutestext for \Page"
\hpword{Máx.}%substitutes text for \Points"
\hsword{Nota}%substitutes text for \Score"
\htword{Total}%substitutes text for \Total:"

\chqword{Pregunta}
\chpword{Máx.}
\chsword{Nota}
\chtword{Total}
\chbpword{Puntos \textit{bonus}}

\renewcommand{\solutiontitle}{\noindent\textbf{Solución:}\par\noindent}


%fuentes
	%\usepackage[T1]{fontenc}
	%\usepackage{fourier}

	\usepackage[T1]{fontenc}
	\usepackage{librebaskerville}


\usepackage{amssymb,amsmath}
\usepackage{graphicx,graphics}
\usepackage{color}
	%	\colorgrids %grids con color
	%	\definecolor{GridColor}{gray}{0.8}
		%grid
%			\setlength{\gridsize}{5mm}
%   			 \setlength{\gridlinewidth}{0.1pt}

%fracciones en texto más elaboradas
	\usepackage{xfrac}

%mejora en la exposición de guarismos
	\usepackage{numprint}
 
%caption
\usepackage[font=small,labelfont=bf,
labelsep=endash]{caption}
	
%soluciones sombreadas	
%\shadedsolutions


\checkboxchar{$\Box$}
\checkedchar{$\blacksquare$}

\usepackage{everypage,draftwatermark} %Marca de agua
\SetWatermarkText{\hspace*{-4cm} \parbox{\textwidth}{\centering  B1\\dist.}}
\SetWatermarkScale{2}
\SetWatermarkLightness{0.9}
%
%\usepackage{background}
%\backgroundsetup{
%scale=1.5,
%angle=180,
%opacity=.05,
%color=black,
%contents={
%\includegraphics[scale=1]{unicorn.jpg}}}


\cfoot{}


%cabecera
\header{\includegraphics[scale=.22]{a.png}}       
      {}
       {$f_2$\textbf{B1-{\small d}}feb24}       
%pie de página

   %enumera en todas menos en la primera
   %\runningfootrule

   %enumera en todas
   %\footrule
   %\footer{}{Página \thepage\ de \numpages}{}
%   \footer{}
%	{Página \thepage\ de \numpages}
%	{\iflastpage{{\footnotesize \textit{Fin del examen}}}{{\footnotesize \textit{Da la vuelta a la hoja}\ldots}}}


	%preguntas test sin indent excesivo	
	\renewcommand{\checkboxeshook}{%
  \settowidth{\leftmargin}{W.}% Included an additional \labelsep + 2.5em skip (now removed)
  \labelwidth\leftmargin\advance\labelwidth-\labelsep%
}


\renewcommand{\questionshook}{\leftmargin=0pt%
  \labelwidth=-\labelsep}
  
  
\begin{document}

\fbox{\parbox{0.945\linewidth}{\textbf{Nombre y apellidos}:\vspace*{20mm}}}

\vspace*{\fill}

\hspace*{1cm}\begin{learn}
{$\triangleright\; M(\chem{O})\simeq16\;g\,mol^{-1}$\newline
$\triangleright\; M(\chem{K})\simeq 39\;g\,mol^{-1}$\newline
$\triangleright\; M(\chem{Cl})\simeq36\;g\,mol^{-1}$\newline
$\triangleright\; M(\chem{N})\simeq14\;g\,mol^{-1},$\newline
$\triangleright\; M(\chem{H})\simeq1\;g\,mol^{-1},$\newline
$\triangleright\; M(\chem{S})\simeq32\;g\,mol^{-1},$\newline
$\triangleright\; M(\chem{Cu})\simeq64\;g\,mol^{-1},$\newline
$\triangleright\; R\simeq 0.082\;atm\,L\,mol^{-1}\,K^{-1}$}
\end{learn}	



\newpage
\vspace*{-1cm}

\begin{problemas}{Un examen es una \textit{comunicación}: debes invocar las leyes en que te basas para resolver los problemas e \textit{indicar brevemente los razonamientos} que sigues para obtener los resultados.}
\end{problemas}
\vspace{1ex}

\qformat{\textbf{Pregunta \thequestion}\quad (\thepoints)\hfill} 

\begin{questions}	
	\question[2]   La descomposición térmica del clorato de potasio se produce según la reacción 
	 $\chem{KClO_3(s)}\rightarrow \chem{KCl(s)}+\chem{O_2(g)}.$
		\begin{parts}
			\part[1] Calcule la entalpía estándar de reacción para la descomposición de un mol de clorato.
			%-22 kJ/mol de clorato. (ajuste 1-> 1 , 3/2)
			\begin{mdframed}[backgroundcolor=yellow!20]  
			El ajuste de la reacción proporciona 
				 $\chem{KClO_3(s)}\rightarrow \chem{KCl(s)}+\frac{3}{2}\chem{O_2(g)},$
			donde hemos elegido los coeficientes para que salga directamente la descomposición de un mol de clorato.
			Con todo la entalpía asociada a la descomposición de un mol de \chem{KClO_3} es
			\[\Delta_r(\chem{KClO_3})=\Delta_fH^0\big(\chem{KCl(s)}\big)-\Delta_fH^0\big(\chem{KClO_3(s)}\big)\simeq -22\;KJ.\]
			\end{mdframed}
			\part[1] Determine la cantidad de calor, a presión constante, que intercambia la reacción 						con el medio al obtenerse $30\;L$ de \chem{O_2} en condiciones estándar de 						presión y temperatura. ¿Se ceden o se absorben?
		%se ceden 17,9 kJ al entorno.
		\begin{mdframed}[backgroundcolor=yellow!20]  
		El calor a presión constante es precisamente la entalpía. En condiciones estándar los $30\;L$ de gas implican que existen $n=\sfrac{PV}{RT}\simeq \frac{1\cdot 30}{0.082\cdot 298}\simeq 1.2\;mol.$
		
		La estequiometría implica que por cada $1.5\,mol$ de \chem{O_2} producido ha debido descomponerse $1\;mol$ de clorato; así que si se producen $1.2\;mol$ de \chem{O_2} se han descompuesto $0.8\;mol$ de \chem{KClO_3}.
		
		Como por cada mol de clorato descompuesto lleva asociada una entalpía estándar de $-22\;KJ$, al descomponerse $0.8\;mol$ el calor a presión constante es de $-17.6\;KJ$, cuyo signo indica que se ceden al entorno (reacción exotérmica).
		\end{mdframed}
		\end{parts}
	\vspace{1ex}
	
		{
		\textbf{Datos:} $\Delta_fH^0\big(\chem{KClO_3(s)}\big)\simeq -414\;kJ\,mol^{-1},\hspace{2ex}							\Delta_fH^0\big(\chem{KCl(s)}\big)\simeq -436\;kJ\,mol^{-1}.$
		}
	\qed

\newpage

	\question[2] Estime la molaridad de:
		\begin{parts}
			\part[1] Una disolución de ácido nítrico etiquetada como al $16\%$ en peso y densidad 								$1.09\;g\,cm^{-3}$.
			\begin{mdframed}[backgroundcolor=yellow!20]  
			La molaridad se calcula como
			\[M\equiv \frac{\text{\small moles de soluto}}{\text{\small litros de disolución}},\]
			de forma que hay que saber cuánto soluto hay en $1\;L$ de disolución. Con el dato de la densidad vemos que  cada litro de disolución tiene una masa de $1090\;g$. De ellos solo el $16\%$ es ácido nítrico, así que hay $174.4\;g$ de \chem{HNO_3}, que equivalen a $\sfrac{174.4}{63}\simeq 2.77\;mol$. 
			
			Así que
			\[M\simeq \frac{2.77\, mol}{1\;L}\simeq 2.8\;mol\,L^{-1}.\]
			\end{mdframed}
			\part[1] La disolución resultante al mezclar $0.5\;L$ de una disolución $2\,M$ de ácido sulfúrico 				con $20\;cm^3$ de $\chem{H_2SO_4}$ concentrado al $96\%$ en masa y densidad $1.84\;g\,cm^{-3}$.
			
			\begin{mdframed}[backgroundcolor=yellow!20]  
			Siguiendo el mismo procedimiento hay que ver cuánto ácido sulfúrico tenemos al mezclas ambas disoluciones. De la primera obtenemos $0.2\;mol\,L^{-1}\cdot 0.5\;L\simeq 1\,mol$. 
			
			En la segunda hay que hacer algún cálculo más: los $20\;cm^3$ de disolución tienen una masa de $20\;cm^3\cdot 1.84\;g\,cm^{-3}\simeq 36.8\;g$. De ellos, el $96\%$ es ácido, así que habrá $0.96\cdot 36.8\,g\simeq 35.3\;g$ de \chem{H_2SO_4}, que equivalen a $\sfrac{35.3}{98}\simeq 0.36\,mol.$
			
			Finalmente
			\[M\simeq\frac{1\,mol+0.36\,mol}{0.5\;L+20\cdot 10^{-3}\,L}\simeq 2.6\;mol\,L^{-1}\].
			\end{mdframed}
		\end{parts}
	\qed

\newpage

	\question[2]   En condiciones normales se hacen reaccionar $60\;g$ de un mineral 		de cobre con una pureza del $90\%$ con $100\;mL$ de una solución $1\,M$ de ácido nítrico. Si el 		rendimiento para la formación de $\chem{NO_2}(g)$ es del $90\%$ y se forma, además, nitrato de 				cobre{\small (II)} y agua:
    	\begin{parts}
    		\part[1] Ajuste la reacción e identifica el reactivo limitante.
    		\begin{mdframed}[backgroundcolor=yellow!20]  
    		La reacción ajustada es
    		\[\chem{Cu+4HNO_3\longrightarrow 2NO_2+Cu(NO_3)_2+2H_2O}.\]
    		Analicemos qué reactivos tenemos: en los $60\;g$ de cobre al $90\%$ hay $54\,g$ puros y, por tanto, $\sfrac{54}{64}\simeq 0.84\;mol$ de \chem{Cu}. Por otro lado en $100\;mL$ de disolucíon $1\;M$ habrá $0.1\;mol$ de \chem{HNO_3}.
    		
    		Como cobre y ácido nítrico reaccionan en relación $1:4$ para hacer reaccionar todo el cobre necesitaríamos $0.84\cdot 4\simeq 3.36\;mol$ de \chem{HNO_3}. 
    		
    		No los tenemos, así que el ácido nítrico es el limitante.
    		\end{mdframed}
    		\part[1] Suponiendo que se comporta idealmente, calcule  el volumen de dióxido de nitrógeno que se 						obtiene.
    		
    		\begin{mdframed}[backgroundcolor=yellow!20]  
    		Al consumirse los $0.1\;mol$ de ácido nítrico se formará la mitad de moles de dióxido de nitrógeno, \textit{i.e.} $\simeq 0.05\;mol$, que darán un volumen de
    		\[V=\frac{nRT}{P}\cdot \frac{90}{100}\simeq \frac{0.05\cdot 0.082\cdot 273}{1}\cdot \frac{90}{100}\simeq 1.0\;L.\]
    		\end{mdframed}
    	\end{parts}
  	\qed

\newpage

	\question[3] El movimiento de un proyectil desde un cierto sistema de referencia colocado en el suelo puede caracterizarse por 	el vector posición dado por
	\[\vec{r}(t)=\big(2+40t,10+30t-5t^2\big)\\;\;\text{\small (SI)}.\]
		\begin{parts}
			\part[1] Calcule la velocidad media durante los primeros $5\;s$ de vuelo.
			\begin{mdframed}[backgroundcolor=yellow!20]  
			La velocidad media es, por definición,
			\[\vec{v}_m\equiv \frac{\vec{r}(t_f)-\vec{r}(t_i)}{t_f-t_i}=\frac{(202,\,35)-(2,\,10)}{5-0}=\frac{(200,\,25)}{5}=(40,5)\;ms^{-1}.\]
			\end{mdframed}
			\part[1] Obtenga las ecuaciones cinemáticas del proyectil.
			\begin{mdframed}[backgroundcolor=yellow!20]  
			Las ecuaciones cinemáticas están formadas, además de por la posición, por el vector velocidad y el vector aceleración, ambos cantidades instantáneas. 			Derivando resulta
			\[\vec{v}(t)=(40,30-10t)\;ms^{-1},\hspace{3ex}\vec{a}(t)=(0,-10)\;ms^{-2}.\]
			\end{mdframed}
			\part[1] ¿A qué altura del suelo está el móvil cuando $t=4\;s$?
			\begin{mdframed}[backgroundcolor=yellow!20]  
			Como el observador está en el suelo la altura en cualquier momento será la componente $y$ de la posición, es decir $10+30t-5t^2$. Para $t=4\;s$ resulta	\(y(t=4\,s)=10+30\cdot 4-5\cdot 4^2=50\;m\).
			\end{mdframed}
		\end{parts}
	\qed

\end{questions}	


\end{document}






